% Generated by scripts/build_latex.py from paper/ks-minimum.md.  Compile with XeLaTeX or LuaLaTeX.
\documentclass[11pt,a4paper]{article}
\usepackage{fontspec}
\setmonofont{DejaVuSansMono.ttf}[Scale=0.82]
\usepackage[margin=1in]{geometry}
\usepackage{amsmath,amssymb}
\usepackage{graphicx}
\usepackage{longtable,booktabs,array,calc}
\usepackage{fvextra}
\usepackage{xurl}
\usepackage[colorlinks=true,linkcolor=black,citecolor=black,urlcolor=blue]{hyperref}
\RecustomVerbatimEnvironment{verbatim}{Verbatim}{breaklines=true,breakanywhere=true,fontsize=\footnotesize}
\providecommand{\tightlist}{\setlength{\itemsep}{0pt}\setlength{\parskip}{0pt}}
\providecommand{\pandocbounded}[1]{#1}
\newcommand{\lt}{<}
\newcommand{\gt}{>}
\usepackage{newunicodechar}
\newunicodechar{μ}{\ensuremath{\mu}}
\newunicodechar{π}{\ensuremath{\pi}}
\newunicodechar{ᵈ}{\ensuremath{^{d}}}
\newunicodechar{⁵}{\ensuremath{^{5}}}
\newunicodechar{⁻}{\ensuremath{^{-}}}
\newunicodechar{¹}{\ensuremath{^{1}}}
\newunicodechar{₀}{\ensuremath{_{0}}}
\newunicodechar{₁}{\ensuremath{_{1}}}
\newunicodechar{₂}{\ensuremath{_{2}}}
\newunicodechar{₃}{\ensuremath{_{3}}}
\newunicodechar{∈}{\ensuremath{\in}}
\newunicodechar{∪}{\ensuremath{\cup}}
\newunicodechar{≤}{\ensuremath{\le}}
\newunicodechar{≥}{\ensuremath{\ge}}
\newunicodechar{⊆}{\ensuremath{\subseteq}}
\newunicodechar{→}{\ensuremath{\to}}
\newunicodechar{∀}{\ensuremath{\forall}}
\newunicodechar{∃}{\ensuremath{\exists}}
\newunicodechar{∧}{\ensuremath{\land}}
\newunicodechar{¬}{\ensuremath{\lnot}}
\newunicodechar{−}{\ensuremath{-}}
\setlength{\emergencystretch}{3em}
\setlength{\parskip}{0.45em}
\setlength{\parindent}{0pt}
\setcounter{secnumdepth}{0}
\renewcommand{\arraystretch}{1.15}
\makeatletter
\def\maxwidth{\ifdim\Gin@nat@width>\linewidth\linewidth\else\Gin@nat@width\fi}
\makeatother
\setkeys{Gin}{width=\maxwidth,keepaspectratio}

\title{Certified exact minima of Kochen-Specker sets inside coordinate pools: 31 and 33 rays in dimension three, 18 in dimension four, 29 in dimension five}
\author{Alexander S. Donahue\\ 0thernes LLC, Washington, USA\\ ORCID 0009-0006-8168-9648\\ \texttt{0\_0@0thernes.art}}
\date{}
\begin{document}
\maketitle
\vspace{-2.5em}\begin{center}\small Version 1.0.0, 10 November 2026. Research preprint, not peer reviewed.\end{center}
\begin{quote}
\textbf{What this paper does and does not do.} The smallest Kochen-Specker (KS) set in dimension three is
an open problem: it has between 24 and 31 rays. This paper does \textbf{not} solve that problem and does
\textbf{not} contain a KS set smaller than a published one. It answers a restricted question exactly:
\emph{inside a given finite pool of rays, what is the smallest KS set?} For the pools of Table 2 the answer
equals the size of the smallest published set of that kind (31, 33, 18 or 29 rays), and the answers
come with certificates that can be re-checked without trusting the search programs; Table 2 says
which kind for each pool. Several of the values were reported before without proof: 31 for integer
coordinates up to 2, 3 and 4, 33 for sixth roots of unity, and 29 as the smallest set generated from
coordinates 0, ±1 in dimension five (Section 1.3). New here, to our knowledge: the proofs with their
certificates, including Lean 4 kernel proofs for two pools; the exact minimum 33 for twelfth,
eighteenth and twenty-fourth roots of unity; and the complete lists of minimum sets for integer
coordinates up to 3 and 4 and for sixth and twelfth roots of unity.

\textbf{How it was made.} The programs, proofs and text were produced by an AI system (Claude, Anthropic)
at the request and under the direction of the author, who is responsible for the release (Section 9).
\end{quote}

\begin{abstract}
A Kochen-Specker set is a finite set of rays admitting no assignment of 0 and 1 that gives exactly one 1
to every orthogonal basis and never 1 to two orthogonal rays. We determine the minimum size of such a
set inside prescribed finite pools of rays. Every KS set must contain a basis missed by any set of
pairwise non-orthogonal rays; this turns the question into a covering problem whose constraints are
individually trivial to verify. Two necessary conditions for critical KS sets, a symmetry normal form,
and a case analysis make it tractable in pools of up to 651 rays. Three independent exact methods are
used and each result carries a certificate. Results: (1) among rays of dimension three with integer
coordinates of absolute value at most 2, 3 or 4, the minimum is 31 and the minimum sets are exactly the
six images of one Conway-Kochen-type set under the symmetries of the cube; (2) among rays with
coordinates 0 or N-th roots of unity, N = 6, 12, 18, 24, the minimum is 33, unique up to symmetry for
N = 6 and 12; (3) among the 40 rays of dimension four with coordinates 0, ±1 the minimum is 18, attained
only by the 16 known sets inside the 24 rays of Peres; (4) among the 121 rays of dimension five with
coordinates 0, ±1 the minimum is 29, the size of the set of Cabello, Estebaranz and García-Alcaine.
The values in (1), the value for N = 6 in (2), and 29 as the smallest set generated from the
coordinates of (4) had been reported without proof. The minimum values in (1) for bound 2 and in (3)
are proved in Lean 4 without Mathlib and checked by its kernel. The remaining statements rest on
unsatisfiability proofs checked by a separate proof checker, with the exceptions listed in Section 4.
Randomised searches in further pools in dimensions three to eight find nothing smaller than the
published sets. The general minimal-KS problem is not settled by any of this.

\end{abstract}
\subsection{1. Introduction}\label{1-introduction}

\subsubsection{1.1 The Kochen-Specker theorem and its minimal sets}\label{11-the-kochen-specker-theorem-and-its-minimal-sets}

Kochen and Specker {[}1{]} and Bell {[}2{]} showed that quantum theory is incompatible with models that assign
every projective measurement a pre-existing outcome independent of the other compatible measurements
performed with it. The combinatorial core is a finite set of rays in a Hilbert space of dimension d ≥ 3
that cannot be 0/1-colored: no assignment gives exactly one 1 to each orthogonal basis while never
giving 1 to two orthogonal rays. Such a set is a \emph{Kochen-Specker set}. The original set had 117 rays in
dimension three; smaller sets followed {[}3, 4, 5{]}, and the smallest KS set in each dimension became a
problem in its own right. KS sets are also the resource behind state-independent contextuality {[}26{]}
and bipartite perfect quantum strategies {[}18, 35{]}.

\subsubsection{1.2 State of knowledge (meta-analysis)}\label{12-state-of-knowledge-meta-analysis}

Table 1 collects the smallest known KS set and the best proven lower bound in each dimension, as far as
we could establish them from the primary sources on 9 and 10 November 2026. Counts follow the standard
definition of Section 2: only the rays of the set are counted. Papers that use hypergraphs with every
basis completed report larger numbers for the same sets in dimension three (the 31-ray set is "51-37").

\textbf{Table 1. Smallest known KS sets and proven lower bounds.}

{\def\LTcaptype{none} % do not increment counter
{\small
\begin{longtable}[]{@{}>{\raggedright\arraybackslash}p{(\linewidth - 8\tabcolsep) * \real{0.0653}}>{\raggedright\arraybackslash}p{(\linewidth - 8\tabcolsep) * \real{0.2302}}>{\raggedright\arraybackslash}p{(\linewidth - 8\tabcolsep) * \real{0.2444}}>{\raggedright\arraybackslash}p{(\linewidth - 8\tabcolsep) * \real{0.2156}}>{\raggedright\arraybackslash}p{(\linewidth - 8\tabcolsep) * \real{0.2444}}@{}}
\toprule\noalign{}
d & smallest known & source & proven lower bound & source \\
\midrule\noalign{}
\endhead
\bottomrule\noalign{}
\endlastfoot
3 & 31 (real) & Conway and Kochen, in {[}3{]} & 24 (real and complex) & {[}10{]}; real case also {[}11{]} \\
3 & 33 with coordinates 0 and roots of unity (14 bases, the fewest known) & {[}20, 21{]} & & \\
4 & 18 & {[}5{]} & 18 & {[}7, 12{]} \\
5 & 29 & {[}6{]} & 18 & {[}12{]} \\
6 & 21 (complex), 31 (real) & {[}13{]}, {[}6{]} & 18 & {[}12{]} \\
7 & 32 (complex), 34 (real) & {[}15{]}, {[}6{]} & 18 & {[}12{]} \\
8 & 34 & {[}14{]} & 18 & {[}12{]} \\
\end{longtable}
}
}

Only dimension four is settled. In dimension three the lower bound rose from 18 {[}8{]} to 22 {[}9{]} to 24
{[}10, 11{]} by exhaustive enumeration of candidate orthogonality graphs, with the final unembeddability
step done by an SMT solver; the authors of {[}11{]} estimate more than a hundred CPU years for the next
order. In dimensions five and above the only proven bound is the dimension-independent 18 of {[}12{]}, and
Pavičić and Waegell {[}32{]} write only that it "seems likely" that the smallest sets in dimensions four to
eight are known.

Restricted results exist in dimension three. Trandafir and Cabello {[}18{]} proved that the Conway-Kochen
set is rigid and that no KS set with 30 or fewer rays arises from the 13-ray set of Yu and Oh {[}26{]} by
their completion procedure, and conjectured that 31 is the minimum. Li, Bright, Trandafir, Cabello and
Ganesh {[}19{]} enumerated, with proof certificates, all KS sets with up to 33 rays containing the 25-ray
completion of the Yu-Oh set. Pavičić and Megill {[}16, 17{]} and Pavičić {[}31{]} generate very large classes
of critical sets from all vectors with components in a small set (their "master" sets) and report the
smallest members they obtain; these are computational reports, without a proof that the generation is
exhaustive and without certificates. Integer vectors can carry KS sets only because their directions
are not rational points of the sphere, which are colorable {[}24, 25{]}; see also {[}23{]}.

\subsubsection{1.3 The pool problem, prior work on it, and what is new here}\label{13-the-pool-problem-prior-work-on-it-and-what-is-new-here}

We fix a finite \emph{pool} of rays and ask for the exact minimum size of a KS set inside it. A pool result
says nothing about rays outside the pool, but the question is exactly solvable, gives matching upper
and lower bounds, and tests the conjectured optimality of the smallest known sets in the places where
they live.

Prior work on this question, as far as we found it:

\begin{itemize}
\tightlist
\item
  Kernaghan {[}22{]}, a preprint whose computations and text were, by its own statement, generated by a
  language model, surveys coordinate alphabets in dimension three. It reports as exact, by a SAT
  procedure for which no certificate is given, the minimum 31 for the 49 rays with coordinates in
  \{0, ±1, ±2\} and the minimum 33 for the 57-ray pool of sixth roots of unity, and it lists 31 also for
  coordinates up to 3. By randomised minimisation it finds six distinct 31-ray sets in the 49-ray pool
  (one remark there speaks instead of 162 distinct sets, six up to automorphism). It proves that the
  pool of N-th roots of unity contains a KS set if and only if 6 divides N. For N = 12, 18, 24 it
  gives only upper bounds, and it has no dimension above three.
\item
  A note on the forum ProblemHive {[}33{]}, posted in October 2026 by an AI agent, reports the minimum 31
  for integer coordinates up to 2, 3 and 4, obtained with CP-SAT and the same kind of coloring cut as
  in Section 3.1. It states itself that its solver logs are not a certificate.
\item
  In dimension four the bound 18 holds for all KS sets {[}7{]}, and in every dimension {[}12{]}. The 18-ray
  sets inside the 24 rays of Peres are known: there are 16 of them {[}39{]}.
\item
  In dimension five, Pavičić and Megill {[}16{]} generate from the components 0, ±1 a master set with 105
  rays and 136 bases, which is the core of our pool B(5, 1), and report that the smallest critical set
  of its class has 29 rays and 16 bases, with two non-isomorphic sets of that size (also {[}32{]}), one
  of which had been obtained in {[}15{]}. No proof that nothing smaller exists is given. Their coloring
  rule constrains only rays that share a basis of the set, so every KS set in their sense is one in
  the sense of Section 2, and a lower bound proved here applies to their sets as well.
\item
  Formal proofs exist for other statements: non-colorability of Peres\textquotesingle s 33 rays in HOL Light {[}36{]},
  the KS paradox in Lean {[}37{]}, case-tree certificates of non-embeddability replayed in Lean for the
  geometric half of the bound 24 {[}34{]}, and a contextuality degree in Lean {[}38{]}. We found no formal
  proof of a minimality statement for KS sets, and none for the Conway-Kochen or the 18-ray set.
\end{itemize}

Against this background we claim as new, to our knowledge and within the search described in
Section 7:

\begin{enumerate}
\def\labelenumi{\arabic{enumi}.}
\tightlist
\item
  proofs with certificates for the results: Lean 4 kernel proofs for the 49-ray pool and for the
  40-ray pool of dimension four, replayable search trees, and unsatisfiability proofs checked by a
  separate checker;
\item
  complete lists of the minimum sets: exactly six for integer coordinates up to 3 and up to 4 (the
  same six as for 2), exactly 3 and 12 for sixth and twelfth roots of unity, each list a single orbit
  under the symmetries of the pool; and the fact that the 40-ray pool of dimension four contains no
  18-ray KS set besides the 16 known ones {[}39{]};
\item
  the exact minimum 33 for twelfth, eighteenth and twenty-fourth roots of unity;
\item
  the proof that 29 is the minimum for coordinates 0, ±1 in dimension five, where 29 had been
  reported as the smallest set obtained {[}16{]}, and the fact that every KS set in that pool contains
  at least three coordinate axes;
\item
  the method beyond the coloring cut itself: necessary conditions for critical sets as constraints,
  a symmetry normal form, an encoding whose unsatisfiability proofs are checked independently, and
  tree certificates with a checker proved sound in Lean (the pattern of certificate and verified
  checker used in {[}34{]}, applied here to a minimality statement).
\end{enumerate}

We do not claim the values 31, 33, 18, 29, which are the sizes of published sets {[}3, 20, 5, 6{]}, nor
the coloring-cut idea, which also appears in {[}33{]} and, in the form of an unsatisfiable-subset loop, in
{[}22{]}.

\subsection{2. Definitions}\label{2-definitions}

A \emph{pool} is a finite set V of rays (one-dimensional subspaces) of a d-dimensional real or complex
inner-product space, d ≥ 3. Two rays are \emph{orthogonal} if their inner product vanishes. A \emph{basis} is a
set of d mutually orthogonal rays of V.

A subset S of V is \emph{colorable} if there is f : S → \{0, 1\} such that

\begin{enumerate}
\def\labelenumi{\arabic{enumi}.}
\tightlist
\item
  f(u) + f(v) ≤ 1 whenever u, v ∈ S are orthogonal, and
\item
  every basis contained in S has a ray with f = 1.
\end{enumerate}

By rule 1 a basis inside S then has exactly one ray of value 1. If no such f exists, S is a
\emph{Kochen-Specker set}.
Only the rays of S count; orthogonal pairs constrain f whether or not they extend to a basis inside S.
Subsets of colorable sets are colorable. A KS set is \emph{critical} if all its proper subsets are colorable;
every KS set contains a critical one, so the minimum is attained by a critical set.

Two families of pools are used. B(d, H) is the set of rays spanned by non-zero integer vectors of \(\mathbb{Z}^d\)
with all coordinates of absolute value at most H. U(d, N) is the set of rays spanned by non-zero vectors
of \(\mathbb{C}^d\) with entries in \(\{0\} \cup \mu_N\), where \(\mu_N\) is the group of N-th roots of unity. Orthogonality is decided
exactly: by integer dot products in B(d, H), and in U(d, N) by reducing the inner product, an integer
combination of N-th roots of unity, modulo the N-th cyclotomic polynomial.

\subsection{3. Method}\label{3-method}

Complete proofs of the statements of this section are in \texttt{paper/PROOFS.md} in the repository; they are
short and are outlined here.

\subsubsection{3.1 Coloring cuts}\label{31-coloring-cuts}

Call I ⊆ V \emph{independent} if no two of its rays are orthogonal, and let Miss(I) be the set of bases of V
containing no ray of I.

\textbf{Lemma 1.} S ⊆ V is a KS set if and only if, for every independent I ⊆ V, S contains a basis of
Miss(I).

If S contains no basis of Miss(I), giving 1 to the rays of I inside S colors S. Conversely the rays of
value 1 of a coloring form an independent set I with no basis of Miss(I) inside S.

So every independent set I yields a \emph{cut}, the clause "S contains at least one basis of Miss(I)", valid
for every KS set and certified by I alone: one checks that I is independent and recomputes Miss(I).
Symmetries of the pool map cuts to cuts. In particular, if C ⊆ V is colorable, every KS set meets the
complement of C.

\subsubsection{3.2 Necessary conditions for critical sets}\label{32-necessary-conditions-for-critical-sets}

\textbf{Lemma 3.} If T is a critical KS set and v ∈ T, then (C1) v lies in a basis contained in T, and (C2)
v is orthogonal to at least d rays of T.

For (C1), a coloring of T without v extends by f(v) = 0 if no basis through v lies in T. For (C2), if v
had at most d − 1 partners they would be the rest of the unique basis b through v; then f(v) = 1 if the
rest of b is all 0, and f(v) = 0 otherwise, extends the coloring.

Applied to the pool itself, the lemma allows the repeated deletion of rays that lie in no basis or have
fewer than d partners; what remains is the \emph{core}. Where the "master" hypergraphs of {[}16, 17{]} exist for
the same vectors, our cores have their sizes: 105 rays and 136 bases for B(5, 1), 1185 and 3596 for
B(5, 2), 332 and 1408 for B(6, 1), 805 and 9936 for B(7, 1), 834 and 1609 for U(6, 3). This is an
external check of the pool construction.

\subsubsection{3.3 Counterexample-guided search}\label{33-counterexample-guided-search}

Take Boolean variables \(x_v\) (ray chosen) and \(y_b\) (basis chosen), with \(y_b\) true exactly when all rays of b
are chosen. For a bound k the \emph{candidate problem} asks for an assignment with at most k rays chosen
that satisfies (C1), (C2) and all cuts collected so far. A solution S is tested: if it is not colorable,
a KS set with at most k rays has been found. Otherwise a coloring of S is extended to an independent
set of the whole pool that misses as few bases as possible; its cut is violated by S and is added,
together with its images under the symmetries of the pool. If the candidate problem becomes infeasible,
the pool contains no KS set with at most k rays, by Lemmas 1 and 3. In a controlled run on B(3, 4) the
conditions (C1) and (C2) reduce the number of rounds from 15 to 6 and the number of cuts from 307 to
123, with no gain in running time (\texttt{results/ablate\_c1c2\_box3\_4\_k30.json}).

\subsubsection{3.4 Search trees}\label{34-search-trees}

For small pools we also use a plain exhaustive search whose trace is a certificate. A node holds an
independent set I and one child for each basis of Miss(I) that can still be added within the budget of
k rays without completing a basis already excluded on the way; the subtree of a child treats the sets
that contain its basis and none of the bases of the earlier siblings. A node without children is closed
because I colors every remaining admissible set. Lemma 5 of \texttt{paper/PROOFS.md} states the checking rule
and proves by induction that an accepted tree shows every set of at most k rays to be colorable. The
rule uses neither (C1) nor (C2). The dual certificate, that one given set is not colorable, is a case
tree on which ray of a basis has value 1.

\subsubsection{3.5 Symmetry and cases}\label{35-symmetry-and-cases}

In B(d, H) the signed coordinate permutations are symmetries. Every set can be brought to a normal form
in which the coordinates are sorted by the number of chosen rays with a non-zero entry there and, for
\(j = 1, \dots, d-1\), the ray \(e_0 - e_j\) is chosen only if \(e_0 + e_j\) is (Lemma 6). For B(5, 1) we also use:

\textbf{Lemma 7.} Every basis of B(5, 1) contains a coordinate axis, and every KS set in B(5, 1) contains at
least three coordinate axes.

The second part holds because the pool with any three axes removed is colorable (an explicit coloring
is in \texttt{results/lemma7\_d5\_axes.json}). A critical KS set therefore has, up to a coordinate permutation,
exactly the axes \(\{e_0, e_1, e_2\}\), or \(\{e_0, e_1, e_2, e_3\}\), or all five, and the three cases are searched
separately. The last case is split further into 81 sub-cases ("cubes") by the states of the four pairs
\((e_0 + e_j,\ e_0 - e_j)\): neither ray chosen, only the first, or both.

\subsubsection{3.6 Three engines and what has to be trusted}\label{36-three-engines-and-what-has-to-be-trusted}

\begin{itemize}
\tightlist
\item
  \textbf{Engine A} solves the candidate problem with the CP-SAT solver of OR-Tools {[}29{]}.
\item
  \textbf{Engine B} encodes the candidate problem from scratch in conjunctive normal form, re-verifying
  each cut from its independent set and using two-sided sequential counters for the counting
  constraints, and solves it with the CDCL solver Glucose {[}28{]} through PySAT {[}27{]}, which emits a DRUP
  proof {[}30{]}. The proof is checked by \texttt{verify/drup\_check.c}, a forward reverse-unit-propagation checker
  of about 300 lines written for this project and tested on correct proofs, corrupted proofs,
  malformed files and satisfiable formulas. Where the register says so, the same formula was also
  given to a second solver, CaDiCaL {[}40{]}, and its proof was checked by the same checker.
\item
  \textbf{Engine C} is the exhaustive tree search of Section 3.4. Its certificates are replayed by an
  independent JavaScript program (\texttt{verify/verify\_cert.mjs}, written from the specification alone,
  which also recomputes orthogonality and all bases from the integer vectors) and, for two pools, by a
  checker whose soundness is proved in Lean 4.
\end{itemize}

For a result marked "proof-checked" the trusted parts are the pool construction, the encoder, the
verification of the cuts, the proof checker, and Lemmas 1, 3, 6, 7; the solvers are not trusted. For a
result marked "Lean" the trusted parts are the Lean kernel and the reading of the theorem statement
and of the few definitions it uses (Section 5); tables, certificates and checkers are not trusted.

\subsection{4. Results}\label{4-results}

\textbf{Table 2. Exact minimum size of a KS set inside each pool.}

{\def\LTcaptype{none} % do not increment counter
{\footnotesize
\begin{longtable}[]{@{}>{\raggedright\arraybackslash}p{(\linewidth - 10\tabcolsep) * \real{0.0968}}>{\raggedright\arraybackslash}p{(\linewidth - 10\tabcolsep) * \real{0.1217}}>{\raggedright\arraybackslash}p{(\linewidth - 10\tabcolsep) * \real{0.0690}}>{\raggedright\arraybackslash}p{(\linewidth - 10\tabcolsep) * \real{0.0879}}>{\raggedright\arraybackslash}p{(\linewidth - 10\tabcolsep) * \real{0.1872}}>{\raggedright\arraybackslash}p{(\linewidth - 10\tabcolsep) * \real{0.4374}}@{}}
\toprule\noalign{}
Pool & rays (core) & bases & minimum & sets of minimum size & evidence \\
\midrule\noalign{}
\endhead
\bottomrule\noalign{}
\endlastfoot
B(3, 2) & 49 (49) & 26 & 31 & 6, one orbit & Lean; tree, 255 nodes; A; B, proof-checked \\
B(3, 3) & 145 (97) & 50 & 31 & 6, one orbit & tree, 11 330 nodes; A; B, proof-checked \\
B(3, 4) & 289 (205) & 110 & 31 & 6, one orbit & A; B, proof-checked \\
U(3, 6) & 57 (57) & 22 & 33 & 3, one orbit & exhaustive tree search; A; B, proof-checked \\
U(3, 12) & 183 (183) & 67 & 33 & 12, one orbit & A; B, proof-checked \\
U(3, 18) & 381 (381) & 136 & 33 & not enumerated & A; B, proof-checked \\
U(3, 24) & 651 (651) & 229 & 33 & not enumerated & A; B with CaDiCaL, proof-checked \\
B(4, 1) & 40 (40) & 32 & 18 & 16, one orbit & Lean; tree, 509 nodes; A; B, proof-checked \\
B(5, 1) & 121 (105) & 136 & 29 & at least 2 classes & B and a second solver, case analysis, proof-checked (Section 4.4) \\
\end{longtable}
}
}

A and B refer to the engines of Section 3.6. The three tree certificates cover the whole pools, so
those rows do not depend on the reduction to the core. For U(3, 24) Glucose ran out of memory and the
formula of engine B was solved by CaDiCaL instead. The counts of minimum sets were obtained by
engine A, excluding each set found together with its orbit and continuing until no candidate was left;
in all six pools the final step (no further KS set of that size) was repeated with engine B and its
proof was checked. The file \texttt{results/RESULTS.md} in the repository is the register behind this table
and records the status of every check.

\subsubsection{4.1 Dimension three, integer coordinates}\label{41-dimension-three-integer-coordinates}

\textbf{Theorem 1.} Let H ∈ \{2, 3, 4\}. Every KS set contained in B(3, H) has at least 31 rays. There are
exactly six KS sets with 31 rays in B(3, H); they have coordinates of absolute value at most 2 and are
the images of one of them under signed coordinate permutations.

The 31-ray sets have 71 orthogonal pairs and 17 bases. The Conway-Kochen set consists of vectors with
coordinates in \{0, ±1, ±2\}, so by the theorem it is one of the six; the count of six agrees with the
six equivalent versions of that set noted in {[}18{]} and with the six sets found by randomised
minimisation in {[}22{]}. For H = 2 the statement "minimum 31" is the Lean theorem
\texttt{KS.Box3H2.minimum\_ks\_size\_vectors}. The minimum 31 for these three pools was reported without
certificates in {[}22{]} (H = 2 as exact, H = 3 in a table) and in {[}33{]} (H = 2, 3, 4); the proofs, and the
statement that for H = 3 and 4 there are no minimum sets besides the six, are new. Enlarging the pool
from 49 to 289 rays adds 156 rays to the core and no new minimum set. The best coloring of the whole
pool misses 1 of 26 bases for H = 2, 2 of 50 for H = 3, 7 of 110 for H = 4 and 14 of 266 for H = 5
(optima of an exact MaxSAT solver, \texttt{results/min\_missed.json}). For H = 5 (577 rays, core 421) the exact search was stopped without an answer after
more than an hour, and we make no claim.

\subsubsection{4.2 Dimension three, roots of unity}\label{42-dimension-three-roots-of-unity}

\textbf{Theorem 2.} For N ∈ \{6, 12, 18, 24\} every KS set contained in U(3, N) has at least 33 rays, and 33 is
attained. In U(3, 6) there are exactly 3 KS sets with 33 rays and in U(3, 12) exactly 12; in both pools
they form a single orbit under coordinate permutations, multiplication of single coordinates by N-th
roots of unity and complex conjugation.

The 33-ray sets have 78 orthogonal pairs and 14 bases, the parameters of the complex set of {[}20, 21{]}.
The value for N = 6 was reported in {[}22{]}; for N = 12, 18, 24 only upper bounds were known. By the
theorem of {[}22{]}, U(3, N) contains a KS set exactly when 6 divides N; our computations agree for
N ≤ 12 (for N = 2, 3, 4, 5, 7, 8, 9, 10, 11 the core is empty). Whether the minimum is 33 for every
multiple of 6 is open.

\subsubsection{4.3 Dimension four}\label{43-dimension-four}

\textbf{Theorem 3.} Every KS set contained in B(4, 1) has at least 18 rays; there are exactly 16 KS sets with
18 rays in B(4, 1), forming one orbit under signed coordinate permutations.

The bound 18 holds for all KS sets in dimension four {[}7{]} and in every dimension {[}12{]}. The 16 sets are
the 16 known 18-ray sets inside the 24 rays of Peres {[}39{]}: none of them uses a ray of the pool with
three non-zero coordinates (\texttt{results/peres24\_check.json}). What the enumeration adds is that these 16
further rays produce no other 18-ray KS set. The theorem is included because the Lean proof of its
first statement, \texttt{KS.Box4H1.minimum\_ks\_size\_vectors}, is to our knowledge the first proof of a
minimality statement for KS sets checked by a proof assistant. The count of 16 sets is not part of
the Lean theorem; it rests on the enumeration and its checked proof (Table 2).

\subsubsection{4.4 Dimension five}\label{44-dimension-five}

\textbf{Theorem 4.} Every KS set contained in B(5, 1) has at least 29 rays, and 29 is attained.

\emph{Proof structure.} By Lemma 7 and the normal form there are three cases for a critical KS set with at
most 28 rays. Table 3 lists how each was closed.

\textbf{Table 3. The three cases of Theorem 4 (bound k = 28).}

{\def\LTcaptype{none} % do not increment counter
{\small
\begin{longtable}[]{@{}>{\raggedright\arraybackslash}p{(\linewidth - 8\tabcolsep) * \real{0.1097}}>{\raggedright\arraybackslash}p{(\linewidth - 8\tabcolsep) * \real{0.1390}}>{\raggedright\arraybackslash}p{(\linewidth - 8\tabcolsep) * \real{0.2122}}>{\raggedright\arraybackslash}p{(\linewidth - 8\tabcolsep) * \real{0.2122}}>{\raggedright\arraybackslash}p{(\linewidth - 8\tabcolsep) * \real{0.3269}}@{}}
\toprule\noalign{}
Case & engine A (CP-SAT) & engine B (Glucose 4.1) & second solver (CaDiCaL 1.9.5) & proof checks \\
\midrule\noalign{}
\endhead
\bottomrule\noalign{}
\endlastfoot
exactly three axes & infeasible, two runs & unsatisfiable, 40 s & unsatisfiable, 28 s & Glucose proof VERIFIED, 94 680 lemmas \\
exactly four axes & stopped without an answer & unsatisfiable, 24 minutes & unsatisfiable, 4 minutes & CaDiCaL proof VERIFIED, 1 614 075 lemmas; Glucose proof VERIFIED, 1 748 269 lemmas \\
all five axes & stopped without an answer & 81 cubes, all unsatisfiable, 84 to 751 s each & 81 cubes, all unsatisfiable, 57 to 449 s each & 81 Glucose proofs: CUBEVERDICT \\
\end{longtable}
}
}

The times are elapsed times of single runs on a loaded machine and vary from run to run; summed over
the 81 cubes they come to about 6 hours for Glucose and 3 hours for CaDiCaL. Both solvers were given
the same formula files, so they do not check the encoding against each other; they check the solvers.

Controls show that the case machinery is not vacuous. With the bound raised to 29, the three-axis
case finds a 29-ray KS set with 16 bases and 151 orthogonal pairs, and the four-axis case finds a
29-ray KS set with 16 bases and 150 orthogonal pairs. These two sets are not isomorphic, in line with
the two smallest critical sets reported for these coordinates in {[}16, 32{]}. The second consists of an
18-ray KS set in a coordinate hyperplane and 11 further rays, the shape of the recursive construction
of {[}6{]}; in the first, no coordinate hyperplane contains a KS set (\texttt{results/d5\_sections.json}).

Critical KS sets with all five axes exist as well, so the third case is not empty: plain random
deletion from the whole pool with the axes deleted last ended with five axes in 3734 of 3828 runs,
with four in 90 and with three in 4 (\texttt{results/d5\_axes\_sanity.json}). The smallest five-axis sets
found have 31 rays, 165 orthogonal pairs and 18 bases; the one recorded in \texttt{results/ks\_sets.json}
contains a KS set of 20 rays and 11 bases in a coordinate hyperplane. As a third control, the search
in one cube of the five-axis case with the bound raised to 31 finds such a set. FIVEAXIS30

Engine A closed only the three-axis case. All three cases were closed by engine B with Glucose and,
on the same formulas, by a second SAT solver, CaDiCaL {[}40{]}; the proof checks are recorded in Table 3
and in the register.

\subsubsection{4.5 Searches for smaller sets}\label{45-searches-for-smaller-sets}

A randomised walk on critical KS sets (remove a few rays, then repair to a critical KS set that keeps
as much of the old set as possible) was run for 10 to 15 minutes in each pool of Table 4. It never
found a set smaller than the smallest published one, and in every pool but one it reached that size,
usually within seconds. The exception is U(3, 24), where the walk stopped at 36 rays although the pool
contains 33-ray sets; this shows how weak a negative result of such a search is.

\textbf{Table 4. Randomised searches (no optimality claim).}

{\def\LTcaptype{none} % do not increment counter
{\small
\begin{longtable}[]{@{}>{\raggedright\arraybackslash}p{(\linewidth - 6\tabcolsep) * \real{0.2663}}>{\raggedright\arraybackslash}p{(\linewidth - 6\tabcolsep) * \real{0.1913}}>{\raggedright\arraybackslash}p{(\linewidth - 6\tabcolsep) * \real{0.2862}}>{\raggedright\arraybackslash}p{(\linewidth - 6\tabcolsep) * \real{0.2562}}@{}}
\toprule\noalign{}
Pool & rays in core & smallest KS set found & smallest published \\
\midrule\noalign{}
\endhead
\bottomrule\noalign{}
\endlastfoot
B(3, 2) & 49 & 31 & 31 {[}3{]} \\
U(3, 12), U(3, 18) & 183, 381 & 33 & 33 {[}20, 21{]} \\
U(3, 24) & 651 & 36 & 33 {[}20, 21{]} \\
B(4, 1), U(4, 6) & 40, 400 & 18 & 18 {[}5{]} \\
B(5, 1), B(5, 2) & 105, 1185 & 29 & 29 {[}6{]} \\
U(5, 4), U(5, 6) & 525, 1505 & 29 & 29 {[}6{]} \\
B(6, 1) & 332 & 31 & 31 real {[}6{]} \\
U(6, 3) & 834 & 21 & 21 {[}13{]} \\
B(7, 1) & 805 & 34 & 34 real {[}6{]} \\
U(7, 3) & 2968 & 32 & 32 {[}15{]} \\
E8 root rays, d = 8 & 120 & 34 & 34 {[}14{]} \\
\end{longtable}
}
}

This is weak evidence, not proof, that these sets are optimal in their pools. We did not decide whether
21 is the minimum in U(6, 3) or 34 in the E8 pool; both exact searches were stopped unfinished. One
further pool was run for which we know no published value: U(5, 3), with a core of 230 rays, where
the walk ended at 85 rays.

\subsection{5. Formal verification in Lean 4}\label{5-formal-verification-in-lean-4}

The Lean development (\texttt{lean/}, Lean 4.35.0, no Mathlib) has four layers.

\begin{enumerate}
\def\labelenumi{\arabic{enumi}.}
\tightlist
\item
  \texttt{KS/Basic.lean} defines pools as tables of bit masks, colorings, and the two certificate checkers
  \texttt{check} (Section 3.4) and \texttt{refute}, and proves their soundness: \texttt{KS.small\_sets\_colorable} (an
  accepted tree certificate implies that every set of at most k rays is colorable) and
  \texttt{KS.not\_colorable\_of\_refute} (an accepted refutation implies that the given set has no coloring).
  While these were being proved, the first versions of both checkers turned out to accept malformed
  certificates; the definitions were corrected, and \texttt{KS/Tests.lean} records that the counterexamples
  are now rejected.
\item
  \texttt{KS/Spec.lean} states the problem without tables. \texttt{boxRays\ d\ H} is the list of integer vectors of
  length d with entries between −H and H whose first non-zero entry is positive and whose entries
  have greatest common divisor 1: one representative of every ray of B(d, H). \texttt{VColorable\ rays\ d\ T}
  says that the set T of positions in the list \texttt{rays} has a coloring, where orthogonal means "dot
  product zero" and a basis is an increasing list of d positions of pairwise orthogonal vectors. Two
  theorems show that the table notion and this vector notion of colorability agree whenever a table
  passes three checks: its orthogonality entries are right, every listed basis consists of d mutually
  orthogonal rays, and every such set is listed.
\item
  Each data module lists the vectors, the orthogonality table and the bases, proves the three checks
  by kernel evaluation, and proves that its list of vectors is \texttt{boxRays\ d\ H} and that every non-zero
  integer vector of the box is parallel to exactly one of them.
\item
  The certificates are evaluated by the kernel (\texttt{decide\ +kernel}). Together with the layers above
  this gives
\end{enumerate}

\begin{verbatim}
theorem KS.Box3H2.minimum_ks_size_vectors :
    (∃ T, T < 2 ^ 49 ∧ card 49 T = 31 ∧ ¬ VColorable (boxRays 3 2) 3 T) ∧
    (∀ T, T < 2 ^ 49 → card 49 T ≤ 30 → VColorable (boxRays 3 2) 3 T)
\end{verbatim}

and the analogous \texttt{KS.Box4H1.minimum\_ks\_size\_vectors} with \texttt{boxRays\ 4\ 1}, 40, 18 and 17. A set of rays
is a natural number T read as a bit mask of positions in the list, and \texttt{card} counts its elements. The
statement mentions no table and no certificate: to trust it one reads the definitions of \texttt{boxRays},
\texttt{dot}, \texttt{card}, \texttt{Sub} and \texttt{VColorable}, about 25 lines of Lean. The axioms used are \texttt{propext},
\texttt{Classical.choice} and \texttt{Quot.sound}; there is no \texttt{sorry} and no \texttt{native\_decide}. Building everything
from scratch takes about one minute. The Lean theorems cover the two minimum values only; the counts
of minimum sets and the other pools rest on the evidence of Table 2.

\subsection{6. An operational check on the Tsotchke stack}\label{6-an-operational-check-on-the-tsotchke-stack}

The 31-ray set was also exercised on four pieces of the Tsotchke software stack, all classical
programs: the Moonlab state-vector simulator (WebAssembly build), the Fubini-Study kernel of the
quantum geometric tensor library, the integer mixers of quantum\_rng as a reproducible pseudo-random
stream, and an exact-integer program in the Eshkol language. A qutrit is embedded in two simulated
qubits. The 71 orthogonal pairs and 17 bases are recovered from simulated overlaps and from
Fubini-Study distances equal to π/2. For 2000 pseudo-random states the 17 basis sums of outcome
probabilities equal 1 to within \(2 \times 10^{-15}\), so the sum over bases is 17 for every state, whereas
exhaustive search shows that an assignment obeying rule 1 gives a 1 to at most 16 of the 17 bases.
This restates the state-independent gap of the set in operational terms. It is a simulation, not an
experiment, and no quantum randomness is involved.

\subsection{7. Limits, prior art, open questions}\label{7-limits-prior-art-open-questions}

\begin{itemize}
\tightlist
\item
  \textbf{Scope.} A pool result excludes smaller KS sets only inside the pool. The general problem in
  dimension three (24 to 31) and in dimensions five and above (18 to the records of Table 1) is open.
\item
  \textbf{Prior art and the word "new".} Literature was searched on 9 and 10 November 2026 on arXiv, Zenodo,
  GitHub and the open web, by topic and by the distinctive numbers. What we found is in Section 1.3.
  Reference {[}33{]} is a forum post that could be read but not independently dated beyond its own
  status line. The first search missed two published facts: the report of 29 as the smallest set for
  the coordinates 0, ±1 in dimension five {[}16{]}, and the 16 sets inside the 24 rays of Peres {[}39{]}. Both
  were found in a review of the manuscript before release, and the claims of Section 1.3 were narrowed
  accordingly. Preprints in this area appear weekly, several of them produced with AI systems, and a
  further overlap we did not find is possible.
\item
  \textbf{What is certified to which level.} Table 2 and the register say which rows have a proof check and
  which rest on solver answers. The randomised searches prove nothing about optimality.
\item
  \textbf{Open questions.} Is 31 the minimum in B(3, H) for all H? Is 33 the minimum in U(3, N) for every
  multiple N of 6? Is 21 the minimum in U(6, 3), and 34 in the E8 pool? Is 29 the minimum in dimension
  five without restriction on the coordinates? Does B(5, 1) contain a KS set with 29 or 30 rays and
  all five coordinate axes?
\end{itemize}

\subsection{8. Reproducibility}\label{8-reproducibility}

Repository: \url{https://github.com/0thernes/kochen-specker-extremal-lab} (site:
\url{https://0thernes.github.io/kochen-specker-extremal-lab/}). The command \texttt{python\ scripts/run\_all.py} runs
the unit tests, regenerates the two small tree certificates and replays them with the JavaScript
verifier, regenerates and checks the unsatisfiability proofs for the two smallest pools, and
re-verifies Lemma 7; \texttt{-\/-lean} adds the Lean build and the axiom audit, \texttt{-\/-stack} the demonstration of
Section 6. Commands and logs for every row of Table 2 are indexed in \texttt{docs/INDEX.md}. Software used:
Python 3.9, PySAT 1.9.dev21 (with Glucose 4.1 and CaDiCaL 1.9.5), OR-Tools 9.15, Node 25, gcc 16,
Lean 4.35.0.

\subsection{9. Disclosure, funding, licence}\label{9-disclosure-funding-licence}

The research, programs, Lean proofs and text were produced by Claude (Anthropic) working as an agent at
the request and under the direction of the author, who takes responsibility for the release. An AI
system is not an author. The work was funded by the author through 0thernes LLC; there was no external
funding and there are no competing interests. The repository is released under the verification
licence in its \texttt{LICENSE} file; third-party components keep their own licences.

\subsection{References}\label{references}

{[}1{]} S. Kochen, E. P. Specker, The problem of hidden variables in quantum mechanics, J. Math. Mech. 17 (1967) 59-87.

{[}2{]} J. S. Bell, On the problem of hidden variables in quantum mechanics, Rev. Mod. Phys. 38 (1966) 447-452.

{[}3{]} A. Peres, Quantum Theory: Concepts and Methods, Kluwer, Dordrecht, 1993 (source of the 31-ray set of J. H. Conway and S. Kochen).

{[}4{]} A. Peres, Two simple proofs of the Kochen-Specker theorem, J. Phys. A 24 (1991) L175-L178.

{[}5{]} A. Cabello, J. M. Estebaranz, G. García-Alcaine, Bell-Kochen-Specker theorem: a proof with 18 vectors, Phys. Lett. A 212 (1996) 183-187.

{[}6{]} A. Cabello, J. M. Estebaranz, G. García-Alcaine, Recursive proof of the Bell-Kochen-Specker theorem in any dimension n \textgreater{} 3, Phys. Lett. A 339 (2005) 425-429; arXiv:quant-ph/0504217.

{[}7{]} M. Pavičić, J.-P. Merlet, B. McKay, N. D. Megill, Kochen-Specker vectors, J. Phys. A 38 (2005) 1577-1592; arXiv:quant-ph/0409014.

{[}8{]} F. Arends, J. Ouaknine, C. W. Wampler, On searching for small Kochen-Specker vector systems (extended version), arXiv:1111.3301 (2011).

{[}9{]} S. Uijlen, B. Westerbaan, A Kochen-Specker system has at least 22 vectors, New Gener. Comput. 34 (2016) 3-23.

{[}10{]} Z. Li, C. Bright, V. Ganesh, A SAT solver + computer algebra attack on the minimum Kochen-Specker problem, Proc. IJCAI 2024, 1898-1906; arXiv:2306.13319.

{[}11{]} M. Kirchweger, T. Peitl, S. Szeider, Co-certificate learning with SAT modulo symmetries, Proc. IJCAI 2023, 1944-1953; arXiv:2306.10427.

{[}12{]} Z.-P. Xu, J.-L. Chen, O. Gühne, Proof of the Peres conjecture for contextuality, Phys. Rev. Lett. 124 (2020) 230401; arXiv:2001.07656.

{[}13{]} P. Lisoněk, P. Badziąg, J. R. Portillo, A. Cabello, Kochen-Specker set with seven contexts, Phys. Rev. A 89 (2014) 042101; arXiv:1308.6012.

{[}14{]} M. Waegell, P. K. Aravind, Parity proofs of the Kochen-Specker theorem based on the Lie algebra E8, J. Phys. A 48 (2015) 225301; arXiv:1502.04350.

{[}15{]} M. Waegell, P. K. Aravind, Minimal complexity of Kochen-Specker sets does not scale with dimension, Phys. Rev. A 95 (2017) 050101; arXiv:1702.05215.

{[}16{]} M. Pavičić, N. D. Megill, Automated generation of arbitrarily many Kochen-Specker and other contextual sets in odd-dimensional Hilbert spaces, Phys. Rev. A 106 (2022) L060203; arXiv:2202.08197.

{[}17{]} M. Pavičić, N. D. Megill, Vector generation of quantum contextual sets in even dimensional Hilbert spaces, Entropy 20 (2018) 928; arXiv:1905.01567.

{[}18{]} S. Trandafir, A. Cabello, Two fundamental solutions to the rigid Kochen-Specker set problem and the solution to the minimal Kochen-Specker set problem under one assumption, Phys. Rev. A 111 (2025) 052204; arXiv:2501.11640.

{[}19{]} Z. Li, C. Bright, S. Trandafir, A. Cabello, V. Ganesh, SAT + NAUTY: orderly generation of small Kochen-Specker sets containing the smallest state-independent contextuality set, arXiv:2604.19947 (2026).

{[}20{]} A. Cabello, Simplest Kochen-Specker set, Phys. Rev. Lett. 135 (2025) 190203; arXiv:2508.07335.

{[}21{]} M. Pavičić, On simplest Kochen-Specker sets, arXiv:2512.10483 (2025).

{[}22{]} M. Kernaghan, The algebraic landscape of Kochen-Specker sets in dimension three, arXiv:2603.16988v9 (2026).

{[}23{]} I. Cortez, C. Morales, M. Reyes, Minimal ring extensions of the integers exhibiting Kochen-Specker contextuality, J. Algebra Appl. (2025), doi:10.1142/S0219498825420174; arXiv:2211.13216.

{[}24{]} C. D. Godsil, J. Zaks, Colouring the sphere, arXiv:1201.0486 (2012).

{[}25{]} D. A. Meyer, Finite precision measurement nullifies the Kochen-Specker theorem, Phys. Rev. Lett. 83 (1999) 3751-3754.

{[}26{]} S. Yu, C. H. Oh, State-independent proof of Kochen-Specker theorem with 13 rays, Phys. Rev. Lett. 108 (2012) 030402.

{[}27{]} A. Ignatiev, A. Morgado, J. Marques-Silva, PySAT: a Python toolkit for prototyping with SAT oracles, Proc. SAT 2018, LNCS 10929, 428-437.

{[}28{]} G. Audemard, L. Simon, Predicting learnt clauses quality in modern SAT solvers, Proc. IJCAI 2009, 399-404.

{[}29{]} L. Perron, F. Didier, CP-SAT, in Google OR-Tools, \url{https://developers.google.com/optimization/cp/cp_solver} .

{[}30{]} N. Wetzler, M. J. H. Heule, W. A. Hunt Jr., DRAT-trim: efficient checking and trimming using expressive clausal proofs, Proc. SAT 2014, LNCS 8561, 422-429.

{[}31{]} M. Pavičić, Quantum contextuality, Quantum 7 (2023) 953; arXiv:2108.00967.

{[}32{]} M. Pavičić, M. Waegell, Generation of Kochen-Specker contextual sets in higher dimensions by dimensional upscaling whose complexity does not scale with dimension and their applications, Phys. Rev. A 110 (2024) 012205; arXiv:2305.08267.

{[}33{]} ProblemHive, post p/xfkay by the account hive-admin, "Smallest Kochen-Specker set in three dimensions: is there one with fewer than 31 vectors?", with claim comments on integer coordinates, October 2026, \url{https://problemhive.org/p/xfkay/thread} (read 10 November 2026).

{[}34{]} S. Siddique, I. Mian, Machine-checked certificates for the geometric half of the minimum Kochen-Specker bound, arXiv:2607.26413 (2026).

{[}35{]} S. Trandafir, A. Cabello, Optimal conversion of Kochen-Specker sets into bipartite perfect quantum strategies, Phys. Rev. A 111 (2025) 022408; arXiv:2410.17470.

{[}36{]} J. Harrison, HOL Light tutorial, file Tutorial/Custom\_tactics.ml (formal proof of the Kochen-Specker paradox for 33 directions), \url{https://github.com/jrh13/hol-light} .

{[}37{]} V. Gupta, FreeWillLean: a formal proof of the Kochen-Specker paradox, \url{https://github.com/guptbot/free_will_lean} (2024).

{[}38{]} Z. Li, A. Gupta, V. Ganesh, The four-qubit hyperbolic quadric has contextuality degree 315, arXiv:2610.26465 (2026).

{[}39{]} M. Waegell, P. K. Aravind, Parity proofs of the Kochen-Specker theorem based on the 24 rays of Peres, Found. Phys. 41 (2011) 1786-1799; arXiv:1103.6058.

{[}40{]} A. Biere, T. Faller, K. Fazekas, M. Fleury, N. Froleyks, F. Pollitt, CaDiCaL 2.0, Proc. CAV 2024, Lecture Notes in Computer Science, Springer, 133-152; doi:10.1007/978-3-031-65627-9\_7.

\end{document}
